% !TEX root=../main.tex
\section{Conclusion and Future Work}\label{sec:conclusion}

Root cause analysis (RCA) is an essential task for OSS operations.
In this work, we formulate RCA as a new causal inference task named \textit{intervention recognition}, based on which, we further obtain the \nameref{thm:rc-criterion} to find the root cause.
We believe such a formulation bridge two well-studied fields (RCA and causal inference) and provide a promising new direction for the critical-yet-hard-to-solve RCA problem in OSS.

% To apply such a criterion in OSS, we propose a novel unsupervised causal inference-based RCA method, \textit{CIRCA}.
To apply such a criterion in OSS, we propose a novel causal inference-based RCA method, \textit{CIRCA}.
CIRCA consists of three techniques, namely structural graph construction, regression-based hypothesis testing, and descendant adjustment.
We verify the theoretical reliability of CIRCA in the simulation study.
Moreover, CIRCA also outperforms baseline methods in a real-world dataset.

In the future, we plan to include faulty data for regression.
We hope that diverse data can help overcome the limited understanding of the system's normal status.
This work rests on a set of assumptions that a real application may not satisfy.
For example, some meta metrics do not have corresponding monitoring metrics in the skeleton we construct for the Oracle database.
% In other words, they are unobserved common drivers of the downstream ones, violating the Causal Sufficiency assumption.
As a result, they can imply common exogenous parents of the downstream monitoring metrics, violating the Markovian assumption.
Explicitly modeling these hidden meta metrics may improve RCA performance.
Meanwhile, the retrospect of analysis mistakes may also point to the lack of monitoring.
Beyond the analysis framework in this work, discoveries on the underlying mechanism of OSS can also help climb the ladder of causation for the RCA task.
